\documentclass[12pt]{article}
\usepackage{amsmath, amssymb}

\title{CSE 400: Fundamentals of Probability in Computing\\
Lecture Scribe — L25: A Tutorial Session – Problem Solving – Inequalities}
\date{April 9, 2026}

\begin{document}

\maketitle

School of Engineering and Applied Science (SEAS), Ahmedabad University

\section*{Q-1}

In a specific city, the average annual income is \$40,000.

\subsection*{(a)}
Use Markov’s inequality to find an upper bound on the probability that a randomly selected person earns \$120,000 or more.

Markov’s inequality:

\[
P(X \ge a) \le \frac{E[X]}{a}
\]

Substituting:

\[
P(X \ge 120,000) \le \frac{40,000}{120,000}
\]

\[
= \frac{1}{3}
\]

\subsection*{(b)}

Can you use Markov’s inequality to bound the probability that a person earns less than \$10,000? Why or why not?

Markov’s inequality applies to probabilities of the form $P(X \ge a)$.

The event $X < 10,000$ is not of this form.

Therefore, Markov’s inequality cannot be used.

\section*{Q-2}

A factory produces resistors with average resistance $\mu = 100\Omega$ and standard deviation $\sigma = 2\Omega$. A resistor is defective if its resistance lies outside $[94\Omega, 106\Omega]$.

\subsection*{(a)}

\[
|X - 100| \ge 6
\]

\[
k = \frac{6}{2} = 3
\]

Chebyshev’s inequality:

\[
P(|X - \mu| \ge k\sigma) \le \frac{1}{k^2}
\]

\[
P(\text{defective}) \le \frac{1}{9}
\]

\subsection*{(b)}

If 10,000 resistors are produced, what is the maximum expected number of defective units?

\[
10000 \cdot \frac{1}{9} = \frac{10000}{9}
\]

\section*{Q-3}

Let $Y$ be a random variable with $0 \le Y \le 1$ and $E[Y] = \mu$. Use the convexity of $g(x) = e^x$ to prove:

\[
E[e^{tY}] \le 1 - \mu + \mu e^t
\]

From convexity:

\[
e^{tY} \le (1 - Y)e^0 + Ye^t
\]

\[
= 1 - Y + Ye^t
\]

Taking expectation:

\[
E[e^{tY}] \le E[1 - Y + Ye^t]
\]

\[
= 1 - E[Y] + E[Y]e^t
\]

\[
= 1 - \mu + \mu e^t
\]

\section*{Q-4}

Let $X > 0$ be a random variable representing server processing time with $E[X] = \mu$. The cost function is $C(X) = X^a$, $a > 1$. Use Jensen’s inequality to prove that:

\[
E[X^a] \ge (E[X])^a
\]

Since $x^a$ is convex for $a > 1$, by Jensen’s inequality:

\[
E[X^a] \ge (E[X])^a
\]

\section*{Q-5}

Show that Chebyshev’s inequality guarantees that:

\[
P(|X - \mu| \ge k\sigma) \le \frac{1}{k^2}
\]

\section*{Q-6}

Suppose we know that the number of items produced in a factory during a week is a random variable with mean 500.

\subsection*{(a)}

Using Markov’s inequality:

\[
P(X \ge 1000) \le \frac{500}{1000} = \frac{1}{2}
\]

\subsection*{(b)}

\[
\sigma = 10
\]

\[
|X - 500| \ge 100
\]

\[
k = \frac{100}{10} = 10
\]

\[
P(|X - 500| \ge 100) \le \frac{1}{100}
\]

\[
P(400 \le X \le 600) \ge 1 - \frac{1}{100} = \frac{99}{100}
\]

\section*{Q-7}

Consider a gambler who, on each play, either wins $+1$ or loses $-1$ unit with equal probability:

\[
P(X_i = 1) = P(X_i = -1) = \frac{1}{2}
\]

Let:

\[
S_n = \sum_{i=1}^{n} X_i
\]

\subsection*{(a)}

\[
E[e^{tX_i}] = \frac{1}{2}e^t + \frac{1}{2}e^{-t}
\]

\[
= \frac{e^t + e^{-t}}{2} \le e^{t^2/2}
\]

\subsection*{(b)}

\[
E[e^{tS_n}] = \prod_{i=1}^{n} E[e^{tX_i}]
\]

\[
\le (e^{t^2/2})^n = e^{nt^2/2}
\]

\section*{Q-8}

A set of 200 people consisting of 100 men and 100 women is randomly divided into 100 pairs of 2 each.

Let $X$ denote the number of pairs consisting of a man and a woman.

An upper bound on:

\[
P(X \le 30)
\]

is obtained using bounds based on deviation from the expected value.

\section*{Q-9}

\[
\frac{1}{1^2 + 1} + \frac{1}{2^2 + 1} + \cdots + \frac{1}{n^2 + 1} \ge \frac{n}{n^2 + 2n + 5}
\]

\section*{Q-10}

Let $X$ be a Poisson random variable with mean 20.

\subsection*{(a)}

\[
P(X \ge 26) \le \frac{20}{26}
\]

\subsection*{(b)}

\[
P(X - 20 \ge 6) \le \frac{20}{36} = \frac{5}{9}
\]

\subsection*{(c)}

\[
P(X \ge 26) \le \inf_{t>0} \frac{E[e^{tX}]}{e^{26t}}
\]

For Poisson:

\[
E[e^{tX}] = e^{20(e^t - 1)}
\]

\[
P(X \ge 26) \le \inf_{t>0} \exp\left(20(e^t - 1) - 26t\right)
\]

\end{document}